
\documentclass{article}
\usepackage{colortbl}
\usepackage{makecell}
\usepackage{multirow}
\usepackage{supertabular}

\begin{document}

\newcounter{utterance}

\centering \large Interaction Transcript for game `referencegame', experiment `line\_grids\_rows', episode 3 with reka{-}flash{-}3.1{-}t0.0.
\vspace{24pt}

{ \footnotesize  \setcounter{utterance}{1}
\setlength{\tabcolsep}{0pt}
\begin{supertabular}{c@{$\;$}|p{.15\linewidth}@{}p{.15\linewidth}p{.15\linewidth}p{.15\linewidth}p{.15\linewidth}p{.15\linewidth}}
    \# & $\;$A & \multicolumn{4}{c}{Game Master} & $\;\:$B\\
    \hline

    \theutterance \stepcounter{utterance}  
    & & \multicolumn{4}{p{0.6\linewidth}}{
        \cellcolor[rgb]{0.9,0.9,0.9}{
            \makecell[{{p{\linewidth}}}]{
                \texttt{\tiny{[P1$\langle$GM]}}
                \texttt{You are given three grids, where each of them is 5 by 5 in size.} \\
\texttt{Grids have empty cells marked with "▢" and filled cells marked with "X".} \\
\texttt{Your task is to generate a referring expression that best describes the target grid while distinguishing it from the two other distractor grids.} \\
\texttt{The first grid is the target grid, and the following two grids are the distractors.} \\
\\ 
\texttt{Target grid:} \\
\\ 
\texttt{□ □ X X X} \\
\texttt{□ □ X X X} \\
\texttt{□ □ X X X} \\
\texttt{□ □ X X X} \\
\texttt{□ □ X X X} \\
\\ 
\texttt{Distractor grid 1:} \\
\\ 
\texttt{□ □ □ X X} \\
\texttt{□ □ □ X X} \\
\texttt{□ □ □ X X} \\
\texttt{□ □ □ X X} \\
\texttt{□ □ □ X X} \\
\\ 
\texttt{Distractor grid 2:} \\
\\ 
\texttt{X □ □ X X} \\
\texttt{X □ □ X X} \\
\texttt{X □ □ X X} \\
\texttt{X □ □ X X} \\
\texttt{X □ □ X X} \\
\\ 
\texttt{Instruction: Describe the target grid.} \\
\texttt{Generate the referring expression starting with the tag "Expression: " for the given target grid. Omit any other text.} \\
            }
        }
    }
    & \\ \\

    \theutterance \stepcounter{utterance}  
    & \multicolumn{4}{p{0.6\linewidth}}{
        \cellcolor[rgb]{0.9,0.9,0.9}{
            \makecell[{{p{\linewidth}}}]{
                \texttt{\tiny{[P1$\rangle$GM]}}
                \texttt{<reasoning>} \\
\texttt{Let's tackle this problem. I need to generate a referring expression that describes the target grid and distinguishes it from the two distractor grids. Let me start by looking at all three grids carefully.} \\
\\ 
\texttt{First, the target grid. All rows are the same: five cells each, with the first two being empty (□) and the last three filled (X). So each row is □ □ X X X. That's five rows of that pattern. So the entire grid is a 5x5 where the left two columns are empty and the right three are filled. So maybe I can describe it as a grid where the left two columns are empty and the right three are filled. But wait, let me check the distractors to make sure that's unique.} \\
\\ 
\texttt{Distractor grid 1: Each row starts with three empty cells, then two filled. So □ □ □ X X. So the first three columns are empty, last two filled. That's different from the target. So the target has two empty on the left, three filled on the right. Distractor 1 has three empty on the left, two filled on the right. So the positions of empty and filled are different.} \\
\\ 
\texttt{Distractor grid 2: Each row starts with an X, then two empty, then two X. So X □ □ X X. So the first column is filled, the next two are empty, then two filled. So the pattern here is different. The target has two empty, three filled. Distractor 2 has one filled, two empty, two filled. So the positions vary.} \\
\\ 
\texttt{So the key difference between the target and the distractors is the number of empty cells on the left and filled on the right. The target has two empty left, three filled right. Distractor 1 has three empty left, two filled right. Distractor 2 has one filled left, two empty, then two filled. So the main distinguishing feature is the number of empty cells in the leftmost columns and the filled cells in the rightmost columns.} \\
\\ 
\texttt{So the referring expression should highlight that the left two columns are empty and the right three are filled. Let me check if that's accurate. Looking at the target grid, each row's first two are empty, last three filled. So columns 1 and 2 are empty, columns 3{-}5 are filled. So the entire grid has} \\
            }
        }
    }
    & & \\ \\

    \theutterance \stepcounter{utterance}  
    & & & \multicolumn{2}{p{0.3\linewidth}}{
        \cellcolor[rgb]{0.9,0.9,0.9}{
            \makecell[{{p{\linewidth}}}]{
                \texttt{\tiny{[GM$|$GM]}}
                \texttt{Invalid generated expression} \\
            }
        }
    }
    & & \\ \\

\end{supertabular}
}

\end{document}
